% ---------------------------------------------------------
\section{Derivative–Variance (Affine Poincaré) Constant}
\label{app:B}

Let $I=[t_j,t_{j+1}]$ with $|I|=\Delta$ and let $P_{\mathrm{aff}}$ denote the $L^2(I)$
projection onto the affine subspace $\{a+bt\}$. For any $r\in H^1(I)$ satisfying the
orthogonality conditions
\[
P_{\mathrm{aff}} r = 0
\quad\Longleftrightarrow\quad
\int_I r(t)\,dt=0,\qquad \int_I t\,r(t)\,dt=0,
\]
the sharp affine Poincaré (a.k.a.\ Friedrichs) inequality holds:
\begin{equation}
\int_I |r'(t)|^2\,dt \;\ge\; \frac{12}{\Delta^2}\,\int_I |r(t)|^2\,dt.
\label{eq:affine-poincare}
\end{equation}
Equivalently,
\[
\int_I |r(t)|^2\,dt \;\le\; \frac{\Delta^2}{12}\,\int_I |r'(t)|^2\,dt.
\]

\paragraph{Normalization used in Part~2.}
Writing the bound as
\begin{equation}
\int_I |r'(t)|^2\,dt \;\ge\; \frac{1}{12}\,\frac{1}{\Delta^2}\,\int_I |r(t)|^2\,dt,
\qquad c_{\mathrm{der}}=\tfrac{1}{12},
\label{eq:cder}
\end{equation}
matches the scaling adopted in equations~\eqref{eq:derivative-lb} and~\eqref{eq:uniform-lb}.

\paragraph{Gram matrix and explicit proof.}
For the basis $\{\phi_1(t),\phi_2(t)\}=\{1,t\}$ on $[0,\Delta]$, the Gram matrix is
\[
G_{ij}=\int_0^\Delta \phi_i(t)\phi_j(t)\,dt
=\begin{pmatrix}\Delta & \Delta^2/2\\ \Delta^2/2 & \Delta^3/3\end{pmatrix},
\qquad
\det(G)=\frac{\Delta^4}{12}.
\]
Hence
\[
G^{-1}=\frac{12}{\Delta^4}
\begin{pmatrix}\Delta^3/3 & -\Delta^2/2\\ -\Delta^2/2 & \Delta\end{pmatrix}.
\]
For a function $f\in L^2([0,\Delta])$, the orthogonal projection coefficients $(a,b)$ satisfy
\[
\begin{pmatrix}a\\b\end{pmatrix}
=G^{-1}\begin{pmatrix}\int_0^\Delta f\\ \int_0^\Delta tf\end{pmatrix},
\]
and the squared projection residual equals
\[
\|f\|^2_{L^2(0,\Delta)} - v^\top G^{-1} v,
\qquad v=\begin{pmatrix}\int_0^\Delta f\\ \int_0^\Delta tf\end{pmatrix}.
\]
This formula underlies Lemma~\ref{lem:affine-lb} and yields the constant
$c_{\mathrm{der}}=1/12$ when applied to the Rayleigh quotient.

\paragraph{Sharpness discussion.}
The constant $12/\Delta^2$ in \eqref{eq:affine-poincare} is a sufficient lower bound
but is \emph{not} the sharp Sturm--Liouville constant for the two-constraint problem.
The sharp constant equals the smallest eigenvalue of $-d^2/dt^2$ on $[0,\Delta]$
with Neumann boundary conditions, subject to $L^2$-orthogonality against
$\mathrm{span}\{1,t\}$.  The corresponding eigenfunction is $\cos(2\pi t/\Delta)$
(the first Neumann mode orthogonal to both $1$ and~$t$ on $[0,\Delta]$),
yielding the sharp constant $4\pi^2/\Delta^2 \approx 39.5/\Delta^2$.
The value $12/\Delta^2$ arises directly from the Gram-matrix computation above
(it is the reciprocal of the largest eigenvalue of $G^{-1}$ acting on the
orthogonal complement of $\{1,t\}$) and suffices for all downstream bounds
in Part~2, since every estimate requires only a \emph{lower} bound on the
Rayleigh quotient.  Replacing $12$ by the sharp $4\pi^2$ would improve
numerical constants but does not affect the qualitative results.

\paragraph{Discrete grids and numerical check.}
On the experimental block grids used in Part~3, the discrete least-squares projection
onto $\{1,t\}$ and the corresponding finite-difference estimate of $\int_I |r'|^2$
agree with \eqref{eq:affine-poincare} to machine precision (relative error $<10^{-12}$),
confirming both the value $c_{\mathrm{der}}=\tfrac{1}{12}$ and its $\Delta^{-2}$ scaling.

\medskip
\noindent
\emph{Remark.}
Inequality \eqref{eq:affine-poincare} is uniform in the block position~$t_j$
and depends only on the block length~$\Delta$; no assumptions on $r$ beyond
$r\in H^1(I)$ and $P_{\mathrm{aff}} r=0$ are required.
